% GENERATED FILE. Edit canonical lessons, then run npm run generate:books.
% canonical-source-hash: fnv1a64:4a0a83b51a0d8a51
% canonical-lessons: RU-C486-zhar, RU-C486-oznob, RU-C486-smorkatsya, RU-C486-okhripnut, RU-C486-pershit

\chapter{A Fever and a Sore Throat}
\label{ch:ru-a-fever-and-a-sore-throat}
\hlchaptermodality{486}

\begin{chapteropening}
\textbf{By the end of this chapter:} \emph{I can describe a fever, chills and a sore throat.}

Complete the last lesson of chapter 486: I can describe a fever, chills and a sore throat.
\end{chapteropening}

\section[\texorpdfstring{\ru{жар} (zhar)}{жар (zhar)}]{\ru{жар} (zhar) — a fever, a high temperature; heat}

\label{lesson:RU-C486-zhar}



\begin{quote}
Before the new one: say the Russian for a faint, then the Russian for to have the shivers.
\end{quote}

\subsection*{You'll want to know: \ru{жар}}
\textbf{\ru{жар}} (\emph{zhar}) — \textquotedblleft{}a fever, a high temperature; heat\textquotedblright{}.

\begin{quote}
\ru{У} \ru{ребёнка} \ru{жар}, \ru{надо} \ru{вызвать} \ru{врача}.
\end{quote}

\begin{quote}
\emph{U rebyónka zhar, nádo výzvat vrachá.}
\end{quote}

\begin{quote}
The child has a fever; we need to call the doctor.
\end{quote}

Masculine. \emph{\ru{У} \ru{него} \ru{жар}} — he has a high temperature.

\begin{cousinweb}
EtymDB lists \emph{\ru{жар}} without a source; it is usually traced to Proto-Slavic \emph{žar\ru{ъ}} 'heat', related to \emph{\ru{гореть}} 'to burn' and \emph{\ru{жарить}} 'to fry'.
\end{cousinweb}

\begin{grammarlens}[title={What you've built}]
One more word for telling, explaining and describing at B1.
\end{grammarlens}

\subsection*{Your turn}
\begin{itemize}
\raggedright
  \item \emph{Say it:} \emph{zhar}
  \item \emph{Say it:} \emph{zhar}, once more
  \item \emph{Say it:} a sentence of your own with \emph{zhar}
  \item \emph{Recall:} say the Russian for a faint, then the Russian for to have the shivers, then say \emph{zhar} again
\end{itemize}

\begin{tcolorbox}[breakable,colback=teal!4,colframe=teal!35!black,title={Before you move on}]
What does \ru{жар} mean? (\textquotedblleft{}A fever, a high temperature; heat\textquotedblright{}.) Say it once more.
\end{tcolorbox}
\section[\texorpdfstring{\ru{озноб} (oznób)}{озноб (oznób)}]{\ru{озноб} (oznób) — chills, shivering}

\label{lesson:RU-C486-oznob}



\begin{quote}
Before the new one: say the Russian for to have the shivers, then the Russian for a fever.
\end{quote}

\subsection*{You'll want to know: \ru{озноб}}
\textbf{\ru{озноб}} (\emph{oznób}) — \textquotedblleft{}chills, shivering\textquotedblright{}.

\begin{quote}
\ru{Вечером} \ru{у} \ru{меня} \ru{начался} \ru{озноб}.
\end{quote}

\begin{quote}
\emph{Vécherom u menyá nachalsyá oznób.}
\end{quote}

\begin{quote}
In the evening I started shivering.
\end{quote}

Masculine.

\begin{cousinweb}
\emph{\ru{озноб}} shares its root with \emph{\ru{знобить}} 'to have the shivers'; EtymDB has no entry, and the origin is not settled.
\end{cousinweb}

\begin{grammarlens}[title={What you've built}]
One more word for telling, explaining and describing at B1.
\end{grammarlens}

\subsection*{Your turn}
\begin{itemize}
\raggedright
  \item \emph{Say it:} \emph{oznób}
  \item \emph{Say it:} \emph{oznób}, once more
  \item \emph{Say it:} a sentence of your own with \emph{oznób}
  \item \emph{Recall:} say the Russian for to have the shivers, then the Russian for a fever, then say \emph{oznób} again
\end{itemize}

\begin{tcolorbox}[breakable,colback=teal!4,colframe=teal!35!black,title={Before you move on}]
What does \ru{озноб} mean? (\textquotedblleft{}Chills, shivering\textquotedblright{}.) Say it once more.
\end{tcolorbox}
\section[\texorpdfstring{\ru{сморкаться} (smorkátsya)}{сморкаться (smorkátsya)}]{\ru{сморкаться} (smorkátsya) — to blow one's nose}

\label{lesson:RU-C486-smorkatsya}



\begin{quote}
Before the new one: say the Russian for a fever, then the Russian for chills.
\end{quote}

\subsection*{You'll want to know: \ru{сморкаться}}
\textbf{\ru{сморкаться}} (\emph{smorkátsya}) — \textquotedblleft{}to blow one's nose\textquotedblright{}.

\begin{quote}
\ru{Он} \ru{громко} \ru{сморкается} \ru{в} \ru{платок}.
\end{quote}

\begin{quote}
\emph{On grómko smorkáyetsya v platók.}
\end{quote}

\begin{quote}
He blows his nose loudly into a handkerchief.
\end{quote}

Imperfective; perfective \ru{высморкаться}.

\begin{cousinweb}
\emph{\ru{сморкаться}} is the reflexive of \emph{\ru{сморкать}} 'to blow one's nose' (EtymDB), the same root as \emph{\ru{насморк}} 'a runny nose'.
\end{cousinweb}

\begin{grammarlens}[title={What you've built}]
One more word for telling, explaining and describing at B1.
\end{grammarlens}

\subsection*{Your turn}
\begin{itemize}
\raggedright
  \item \emph{Say it:} \emph{smorkátsya}
  \item \emph{Say it:} \emph{smorkátsya}, once more
  \item \emph{Say it:} a sentence of your own with \emph{smorkátsya}
  \item \emph{Recall:} say the Russian for a fever, then the Russian for chills, then say \emph{smorkátsya} again
\end{itemize}

\begin{tcolorbox}[breakable,colback=teal!4,colframe=teal!35!black,title={Before you move on}]
What does \ru{сморкаться} mean? (\textquotedblleft{}To blow one's nose\textquotedblright{}.) Say it once more.
\end{tcolorbox}
\section[\texorpdfstring{\ru{охрипнуть} (okhrípnut)}{охрипнуть (okhrípnut)}]{\ru{охрипнуть} (okhrípnut) — to go hoarse}

\label{lesson:RU-C486-okhripnut}



\begin{quote}
Before the new one: say the Russian for chills, then the Russian for to blow one's nose.
\end{quote}

\subsection*{You'll want to know: \ru{охрипнуть}}
\textbf{\ru{охрипнуть}} (\emph{okhrípnut}) — \textquotedblleft{}to go hoarse\textquotedblright{}.

\begin{quote}
\ru{После} \ru{концерта} \ru{певец} \ru{охрип}.
\end{quote}

\begin{quote}
\emph{Pósle kontsérta pevéts okhríp.}
\end{quote}

\begin{quote}
After the concert the singer went hoarse.
\end{quote}

Perfective; past \ru{охрип}, \ru{охрипла}. Imperfective \ru{хрипнуть}.

\begin{cousinweb}
\emph{\ru{охрипнуть}} is built on \emph{\ru{о}-} and \emph{\ru{хрипнуть}}, from \emph{\ru{хрип}} 'a wheeze'; EtymDB lists \emph{\ru{хрип}} and \emph{\ru{хрипнуть}} without a source, and the root is usually taken as imitative.
\end{cousinweb}

\begin{grammarlens}[title={What you've built}]
One more word for telling, explaining and describing at B1.
\end{grammarlens}

\subsection*{Your turn}
\begin{itemize}
\raggedright
  \item \emph{Say it:} \emph{okhrípnut}
  \item \emph{Say it:} \emph{okhrípnut}, once more
  \item \emph{Say it:} a sentence of your own with \emph{okhrípnut}
  \item \emph{Recall:} say the Russian for chills, then the Russian for to blow one's nose, then say \emph{okhrípnut} again
\end{itemize}

\begin{tcolorbox}[breakable,colback=teal!4,colframe=teal!35!black,title={Before you move on}]
What does \ru{охрипнуть} mean? (\textquotedblleft{}To go hoarse\textquotedblright{}.) Say it once more.
\end{tcolorbox}
\section[\texorpdfstring{\ru{першить} (pershít)}{першить (pershít)}]{\ru{першить} (pershít) — to tickle (in the throat)}

\label{lesson:RU-C486-pershit}



\begin{quote}
Before the new one: say the Russian for to blow one's nose, then the Russian for to go hoarse.
\end{quote}

\subsection*{You'll want to know: \ru{першить}}
\textbf{\ru{першить}} (\emph{pershít}) — \textquotedblleft{}to tickle (in the throat)\textquotedblright{}.

\begin{quote}
\ru{У} \ru{меня} \ru{в} \ru{горле} \ru{першит}.
\end{quote}

\begin{quote}
\emph{U menyá v górle pershít.}
\end{quote}

\begin{quote}
I've got a tickle in my throat.
\end{quote}

Impersonal: \emph{\ru{в} \ru{горле} \ru{першит}}. Imperfective only.

\begin{cousinweb}
\emph{\ru{першить}} is usually linked to \emph{\ru{перхать}} 'to clear one's throat'; EtymDB has no entry, so the origin is not settled.
\end{cousinweb}

\begin{grammarlens}[title={What you've built}]
One more word for telling, explaining and describing at B1.
\end{grammarlens}

\subsection*{Your turn}
\begin{itemize}
\raggedright
  \item \emph{Say it:} \emph{pershít}
  \item \emph{Say it:} \emph{pershít}, once more
  \item \emph{Say it:} a sentence of your own with \emph{pershít}
  \item \emph{Recall:} say the Russian for to blow one's nose, then the Russian for to go hoarse, then say \emph{pershít} again
\end{itemize}

\begin{tcolorbox}[breakable,colback=teal!4,colframe=teal!35!black,title={Before you move on}]
What does \ru{першить} mean? (\textquotedblleft{}To tickle (in the throat)\textquotedblright{}.) Say it once more.
\end{tcolorbox}
